% CE6190 Group Project 2 -- REPORT TEMPLATE, not a completed experiment report.
% Build instructions are in README.md. All measured values are intentionally absent.
% Keep the draft labels until every claim and table has been replaced by own evidence.
\documentclass[11pt,a4paper]{article}
\usepackage[margin=24mm]{geometry}
\usepackage[T1]{fontenc}
\usepackage{lmodern}
\usepackage{microtype}
\usepackage{amsmath,amssymb}
\usepackage{booktabs,tabularx,array,longtable}
\usepackage{graphicx}
\usepackage{xcolor}
\usepackage{enumitem}
\usepackage{fancyhdr}
\usepackage{natbib}
\usepackage[hidelinks]{hyperref}
\usepackage{url}
\definecolor{draftorange}{RGB}{143,68,0}
\definecolor{draftbg}{RGB}{255,246,231}
\newcommand{\pending}{\textcolor{draftorange}{\textsf{TBD}}}
\newcommand{\fillin}[1]{\textcolor{draftorange}{[\textsf{Fill: #1}]}}
\newcommand{\instruction}[1]{\par\smallskip\noindent
  \fcolorbox{draftorange}{draftbg}{\parbox{\dimexpr\linewidth-2\fboxsep-2\fboxrule\relax}{
  \small\textbf{Template instruction.} #1}}\par\smallskip}
\newcommand{\placeholderfigure}[2]{\begin{figure}[htbp]\centering
  \fbox{\parbox[c][32mm][c]{0.92\linewidth}{\centering\textbf{Pending own measurements}\par\medskip#1}}
  \caption{#2}\end{figure}}
\newcolumntype{Y}{>{\raggedright\arraybackslash}X}
\setlength{\parindent}{0pt}
\setlength{\parskip}{5pt}
\setlength{\headheight}{15pt}
\setlist{nosep,leftmargin=*}
\pagestyle{fancy}
\fancyhf{}
\fancyhead[L]{\small CE6190 Group Project 2}
\fancyhead[R]{\small\textcolor{draftorange}{TEMPLATE: RESULTS PENDING}}
\fancyfoot[C]{\thepage}
\title{Efficient Text-to-Video Generation via Adaptive Caching\\
\large Comparative Study of SeaCache and SenCache\\[4pt]
\textcolor{draftorange}{\textbf{REPORT TEMPLATE -- NO GPU RESULTS}}}
\author{Group \fillin{XX}\\\fillin{actual member names and student identifiers}}
\date{CE6190, 2026 -- prepared for completion after experiments}

\begin{document}
\maketitle
\thispagestyle{fancy}
\instruction{This is a writing and analysis scaffold. No GPU experiment results, measured
speedups, or quality scores are provided. Do not submit it as a completed project report.
Replace each TBD, instruction box, and figure placeholder with verified results from your
own runs. Do not replace missing measurements with numbers copied from a paper.}

\begin{abstract}
This template specifies a comparative study of adaptive caching for text-to-video
generation with a shared pretrained Wan checkpoint. The planned comparison includes
uncached inference, TeaCache, MagCache, SeaCache, and SenCache. It separates agreement
with uncached outputs from video quality and text alignment, and distinguishes
calibration cost from inference latency. The proposed analysis examines matched-budget
quality, hyperparameter sensitivity, key-component ablations, and failure cases.
\textbf{No performance conclusion is asserted here: experiments and measurements are pending.}
\end{abstract}
\instruction{Rewrite the abstract last. State the exact model, benchmark, experimental
scope, measured findings with uncertainty where appropriate, and the main limitation.}

\section{Introduction}
Video generation repeatedly evaluates a large network while updating a video latent.
Caching can reduce this cost by reusing intermediate computations at selected denoising
steps. The central decision is whether a previously computed result remains an adequate
approximation for the current state. SeaCache uses a representation aligned with the
spectral evolution of denoising~\citep{chung2026seacache}; SenCache estimates change using
input drift and calibrated model sensitivity~\citep{haghighi2026sencache}.

This study is scoped to efficient text-to-video generation. It does not change the
generation task to autoregressive video prediction, and it does not overlap with the
individual project on in-context segmentation. The following research questions are
\emph{proposed questions}, not already established contributions:
\begin{enumerate}[label=RQ\arabic*.]
\item At a matched measured latency, how do the caching methods differ in fidelity to
uncached outputs, video quality, text alignment, and motion?
\item Which decision components explain the observed trade-offs, after accounting for
refresh safeguards, CFG handling, and calibration?
\item When do sample-dependent schedules help, and which content or generation settings
cause degradation or unreliable error estimates?
\end{enumerate}
\instruction{Replace or narrow the questions to those actually answered. Write the final
contributions as supported findings, not promises to outperform previous methods.}

\section{Background and Methods}
\subsection{Generation, network computation, and cache reuse}
Let $z_t$ be the video latent at noise time $t$, $c$ the fixed text condition, and $\theta$
the pretrained parameters. Wan uses a flow-matching formulation~\citep{wan2025wan}; a
network prediction and a sampler update can be written abstractly as
\begin{equation}
  p_t=v_\theta(z_t,t,c),\qquad z_s=\operatorname{Update}(z_t,p_t,t,s),\quad s<t.
\end{equation}
The precise update depends on the chosen sampler. The denoising-time index is distinct
from the video-frame index. Caching may skip expensive internal network work while
the latent continues to update at the selected sampling steps. Inference does not
update $\theta$.

For a cached block stack, let $h_k^{\mathrm{in}}$ and $h_k^{\mathrm{out}}$ be its input and
output at denoising iteration $k$. A residual-cache abstraction is
\begin{equation}
  r_k=h_k^{\mathrm{out}}-h_k^{\mathrm{in}},\qquad
  \widetilde h_k^{\mathrm{out}}=h_k^{\mathrm{in}}+r_j,\quad j<k,
\end{equation}
where $j$ is the most recent refresh iteration. This preserves the new input while
reusing an old residual. The exact cache boundary must be verified in the selected
implementation. A pedagogical example of reusing the final noise prediction must not
be presented as the actual Wan implementation.

\subsection{Compared policies}
\textbf{Uncached reference.} Evaluate the full block stack at every selected denoising
step, with the same checkpoint and generation settings as the cached runs.

\textbf{TeaCache.} The paper uses inexpensive input-side information with timestep
conditioning and a rescaling strategy to estimate output differences~\citep{liu2025teacache}.
The exact proxy is implementation-dependent: document whether the chosen Wan path reads
a latent-dependent feature or the timestep embedding/projection. Do not infer
sample-adaptivity merely from the method name.

\textbf{MagCache.} Magnitude-ratio information is used to model residual change and
guide reuse~\citep{ma2025magcache}. Record the calibration source, ratio table, skip
limit, and any stage-specific protection in the selected implementation.

\textbf{SeaCache.} A timestep-aligned spectral filter is applied to an input-feature
representation before measuring change~\citep{chung2026seacache}. In the audited Wan
path, successive filtered-feature differences contribute to a cumulative refresh
score. The filter is a decision aid; it is not an instruction to low-pass-filter the
final generated video. Record the actual feature location, normalization, schedule,
threshold, and reset behavior.

\textbf{SenCache.} The decision combines latent drift and timestep drift with calibrated
sensitivity profiles~\citep{haghighi2026sencache}. A schematic score at current time $t$
relative to the last refresh time $t_r$ is
\begin{equation}
 E_{\mathrm{Sen}}=s_x(t_r)\lVert z_t-z_{t_r}\rVert_2+s_t(t_r)|t-t_r|.
\end{equation}
The scalar profiles summarize model response to perturbations. The practical finite
difference estimates are not exact global Jacobian operator norms, and the decision
threshold is not a guarantee on final video error. Calibration prepares profiles;
it does not retrain the generator.

\instruction{Cross-check this section against the exact code actually run. Add the
implementation revisions, CFG branch rule, first/last-step protections, reset behavior,
and strict versus non-strict threshold comparison. Label all modified policies as
variants. Consult the accompanying study guide for the detailed equation derivations.}

\section{Experimental Protocol}
\subsection{Data separation and shared settings}
VBench provides prompts and dimension-specific evaluation protocols~\citep{huang2024vbench}.
The proposed main test uses a fixed VBench prompt manifest. Calibration data, validation
prompts for threshold selection, and final evaluation prompts must be separated.
If all VBench prompts are reserved for testing, choose an independent validation source.
Do not assume that two papers using VBench used identical prompt counts or settings.

\begin{table}[htbp]
\centering\small
\begin{tabularx}{\linewidth}{@{}p{0.30\linewidth}Y@{}}\toprule
Field & Actual setting to report \\\midrule
Checkpoint and revision/hash & \fillin{exact model ID, weights revision, checksum}\\
Implementation revisions & \fillin{Wan and each caching repository commit; local changes}\\
Prompts and split & \fillin{metadata revision, unique prompt count, prompt IDs, split seed}\\
Seeds per prompt & \fillin{complete seed list and initial-noise pairing rule}\\
Resolution, frames, FPS & \fillin{actual generation and saved-video settings}\\
Sampling & \fillin{sampler, steps, timestep schedule/shift, CFG, negative prompt}\\
Numerics and memory settings & \fillin{precision, attention backend, offload, compilation}\\
Hardware and software & \fillin{GPU model/count, driver, CUDA, PyTorch, dependencies}\\
Calibration & \fillin{data IDs, sample count, profile/table version, runtime}\\
Validation selection & \fillin{threshold grid, equal search budget, selection objective}\\
Run accounting & \fillin{successful/failed runs, retries, exclusion rule}\\\bottomrule
\end{tabularx}
\caption{Shared settings: pending actual configuration. A method-specific difference
must be explained rather than silently included in this shared table.}
\label{tab:protocol}
\end{table}

\subsection{Metrics and timing}
Report both \textbf{agreement with the uncached reference} (PSNR, SSIM, and LPIPS under
a fixed frame-alignment and preprocessing protocol) and \textbf{video generation
quality} (the official VBench aggregation and relevant dimensions). Agreement does not
establish text correctness, realism, or good motion. High temporal consistency alone
can reward a nearly static video, so inspect dynamic degree and actual motion jointly.

State the PSNR data range, SSIM implementation, LPIPS network/version, resizing, frame
sampling, and whether metrics are computed before or after video compression. Use
matching frames from matching prompt/seed pairs. Preserve the official VBench
dimension-to-prompt mappings; a custom subset score must not be labeled a standard
full-benchmark score.

Measure GPU-synchronized wall time with warm-up and repetitions. Define whether each
reported time includes text encoding, denoising, VAE decoding, video encoding, loading,
and device transfers. Retain per-run measurements, and use the same hardware and
timing boundary for all comparisons. Record sampling steps, full block-stack branch
evaluations, reuse calls, and decision overhead separately; CFG may require two
branches per sampling step.
\begin{equation}
\operatorname{Speedup}(m)=\frac{\overline T_{\mathrm{uncached}}}{\overline T_m},
\qquad
T_{\mathrm{amortized}}(m,N)=\overline T_m+\frac{T_{\mathrm{calibration},m}}{N}.
\end{equation}
The bars above denote the predeclared aggregate of paired-run latency measurements.
State the aggregation explicitly; the ratio of means is not the mean of per-sample
speedups. $N$ is the number of generated videos sharing the same calibration.

\subsection{Matched budgets and statistical analysis}
Tune each method on the same validation budget; thresholds have different units and
must not be forced to the same numerical value. Freeze selected settings before
testing. Plot the complete evaluated latency--quality trade-off rather than choosing
only a favorable threshold. Report achieved budgets, not just requested targets.

Use paired prompt/seed comparisons. For uncertainty, resample at the prompt level
when several seeds share a prompt; do not treat individual video frames as independent
experimental samples. Report the chosen estimator, confidence interval procedure,
sample count, and missing-run treatment. If a small pilot is all that has been run,
label it a pilot and restrict the conclusion accordingly.

\begin{table}[htbp]\centering\small
\begin{tabularx}{\linewidth}{@{}lYYYY@{}}\toprule
Budget & Target latency & Selection split & Achieved latency & Quality outcome\\\midrule
A & \pending & \pending & \pending & \pending\\
B & \pending & \pending & \pending & \pending\\
C & \pending & \pending & \pending & \pending\\\bottomrule
\end{tabularx}
\caption{Budget plan, not results. Populate a row per method within each budget after
validation selection and held-out evaluation. Remove unused budget levels.}
\end{table}

\section{Results and Analysis: Pending Own Measurements}
\instruction{All tables below are empty result schemas. Insert own measurements only.
Each row must identify its configuration and actual sample count. Distinguish
implementation validation, pilot runs, and the final test. Do not write winner claims
until uncertainty and matched-budget comparisons have been checked.}

\subsection{Efficiency and calibration cost}
\begin{table}[htbp]\centering\small
\begin{tabularx}{\linewidth}{@{}lYYYYY@{}}\toprule
Method & Config ID & Time (s) & Speedup & Full-stack calls & Calib. (s)\\\midrule
Uncached & \pending & \pending & reference & \pending & n/a\\
TeaCache & \pending & \pending & \pending & \pending & \pending\\
MagCache & \pending & \pending & \pending & \pending & \pending\\
SeaCache & \pending & \pending & \pending & \pending & \pending\\
SenCache & \pending & \pending & \pending & \pending & \pending\\\bottomrule
\end{tabularx}
\caption{Efficiency schema. No values have been measured. Report the latency
aggregation and uncertainty; specify whether call counts include both CFG branches.
Explain n/a, reused published calibration, and locally measured calibration separately.}
\end{table}

\subsection{Fidelity and video quality}
\begin{table}[htbp]\centering\small
\begin{tabularx}{\linewidth}{@{}lYYYY@{}}\toprule
Method & Paired videos & PSNR $\uparrow$ & SSIM $\uparrow$ & LPIPS $\downarrow$\\\midrule
TeaCache & \pending & \pending & \pending & \pending\\
MagCache & \pending & \pending & \pending & \pending\\
SeaCache & \pending & \pending & \pending & \pending\\
SenCache & \pending & \pending & \pending & \pending\\\bottomrule
\end{tabularx}
\caption{Agreement with the matching uncached output: pending measurements. The
uncached output is the reference, not ground-truth video. State the unit of averaging.}
\end{table}

\begin{table}[htbp]\centering\small
\begin{tabularx}{\linewidth}{@{}lYYYYY@{}}\toprule
Method & Total & Quality & Semantic & Motion smoothness & Dynamic degree\\\midrule
Uncached & \pending & \pending & \pending & \pending & \pending\\
TeaCache & \pending & \pending & \pending & \pending & \pending\\
MagCache & \pending & \pending & \pending & \pending & \pending\\
SeaCache & \pending & \pending & \pending & \pending & \pending\\
SenCache & \pending & \pending & \pending & \pending & \pending\\\bottomrule
\end{tabularx}
\caption{VBench schema: pending measurements. Add all other evaluated dimensions in
the supplement. State whether scores use $[0,1]$ or percentage units. Only use the
official aggregate when its required dimensions and protocol have been completed.}
\end{table}

\placeholderfigure{Insert measured latency--quality curves, with one point per frozen
configuration and an uncached reference. Identify the evaluated range and uncertainty.}
{Placeholder for the main trade-off figure. This box is not an experimental plot.}

\instruction{Discuss whether rankings agree across fidelity, video quality, and motion;
whether differences exceed run-to-run variation; and where a speedup is practically
meaningful. Trace every statement to a table, figure, or logged case. Negative or mixed
results are valid results.}

\subsection{Hyperparameter analysis and component ablations}
The proposed ablations should answer specific causal questions. Examples include
removing SeaCache's spectral filtering, changing its timestep dependence, removing one
SenCache sensitivity term, and varying calibration sample count. These changes are
\emph{proposed variants}; their availability and results are not asserted here. Match
latency or full-stack evaluation budget and document all altered safeguards.

\begin{table}[htbp]\centering\small
\begin{tabularx}{\linewidth}{@{}p{0.22\linewidth}YYYY@{}}\toprule
Variant & Question tested & Budget control & Main metric & Outcome\\\midrule
Sea: no filter & \fillin{specific hypothesis} & \pending & \pending & \pending\\
Sea: fixed filter & \fillin{specific hypothesis} & \pending & \pending & \pending\\
Sen: latent term only & \fillin{specific hypothesis} & \pending & \pending & \pending\\
Sen: time term only & \fillin{specific hypothesis} & \pending & \pending & \pending\\
Sen: calibration size & \fillin{specific hypothesis} & \pending & \pending & \pending\\
Additional variant & \fillin{actual implemented change} & \pending & \pending & \pending\\\bottomrule
\end{tabularx}
\caption{Ablation plan: remove unimplemented rows, add the actual configuration, and
replace every outcome with measured evidence before submission.}
\end{table}

Potential diagnostics include prompt-by-step refresh matrices, decision-score traces,
and local approximation errors on a small diagnostic set. An extra full evaluation at
a skipped step may provide an offline error target, but its cost must not be hidden in
the timed accelerated run. Keep diagnostic runs and latency runs separate.

\subsection{Qualitative and failure-case analysis}
\placeholderfigure{Insert synchronized frames from the same prompt and seed for all
methods. Include an independently selected success case and failure case, and identify
frame numbers, thresholds, elapsed time, and any crops.}
{Placeholder for own generated videos. A running dog is a suitable example; preserve
the uncropped context and inspect motion as well as static appearance.}

\instruction{Describe the actual error: identity drift, blurred texture, wrong object,
flicker, frozen motion, motion discontinuity, or text mismatch. Relate it to the logged
refresh behavior without claiming causation from a single example. Explain how cases
were selected and retain unsuccessful cases.}

\section{Discussion, Limitations, and Future Extension}
\fillin{Summarize the supported advantages and disadvantages of each method at the
tested budgets. Separate the paper's intended mechanism, the selected repository's
behavior, and your observed results.}

Limitations to assess include coverage of prompts and motion types, a single backbone
or checkpoint size, calibration transfer across frame counts or schedulers, metric
limitations, statistical uncertainty, and implementation overhead. Report unexpected
failures and any material discrepancy from official code. Resource availability does
not remove the need to define the scope of a claim.

One possible extension is to test whether frequency-based and sensitivity-based signals
provide complementary information. If implemented, define the combined rule, keep
independent state/reset logic correct, include its overhead, and compare with both
individual policies at matched latency. Present this as a project experiment unless
novelty has been established by a separate literature review. A negative result should
be analyzed rather than omitted.

\section{Conclusion}
\fillin{Answer RQ1--RQ3 using only completed measurements. State the main trade-off,
which settings it holds under, and the most important remaining limitation. No result
or winner is currently available in this template.}

\section*{Team Contributions}
\instruction{Enter the actual members and verifiable contributions. Do not invent a
group size or use equal-contribution statements without agreement. All members should
understand the task, methods, and conclusions.}
\begin{tabularx}{\linewidth}{@{}p{0.26\linewidth}YY@{}}\toprule
Member & Contribution & Cross-check / evidence\\\midrule
\fillin{name} & \fillin{completed work} & \fillin{review or commit/run record}\\
\fillin{add actual members} & \fillin{completed work} & \fillin{review or record}\\\bottomrule
\end{tabularx}

\section*{Reproducibility Statement}
\fillin{List the code/config revision, environment lock, prompt manifest hash, raw-run
records, calibration provenance, and scripts used to regenerate each reported table.
State what was successfully reproduced in a clean environment.}

\bibliographystyle{plainnat}
\bibliography{references}

\appendix
\section{Supplementary Evidence to Add}
\begin{itemize}
\item Complete per-dimension metrics, per-prompt aggregates, confidence intervals, and
the method used to compute them.
\item All evaluated thresholds, calibration settings, exact protected steps, skip
limits, and CFG branch rules.
\item Source revisions and local patch descriptions; comparison to the uncached path.
\item Failure/retry manifest, exclusion rules, and additional qualitative examples.
\item Mapping from each reported table/figure to its command and input result files.
\end{itemize}
\instruction{Delete empty supplementary sections. The course requires a report of at
least four pages, with introduction, methods, experiments/analysis, and conclusion.
Template pages and empty placeholders are not completed analysis. Final submission:
\texttt{group\_XX\_report\_project2.pdf} and \texttt{group\_XX\_code\_project2.zip}.
Submit actual files; do not substitute cloud or GitHub links. Do not include datasets
or model weights in the code ZIP.}
\end{document}
